%% ============================================================
%% TSA MEETING — DETAILED SPEAKER NOTES
%% Presentation: Self-Adaptive Model-Based Protection for
%%               Future Power Systems
%% Author:       Anoop V Eluvathingal
%% IIT Madras, Dept. Electrical Engineering
%% June 2026
%% ============================================================

\documentclass[11pt,a4paper]{article}

\usepackage{geometry}
\geometry{top=25mm, bottom=25mm, left=28mm, right=28mm}

\usepackage{amsmath,amssymb}
\usepackage{booktabs}
\usepackage{array}
\usepackage{xcolor}
\usepackage{graphicx}
\usepackage{titlesec}
\usepackage{parskip}
\usepackage{microtype}
\usepackage{fancyhdr}
\usepackage{enumitem}
\usepackage{mdframed}
\usepackage{hyperref}
\hypersetup{colorlinks=true, linkcolor=blue!60!black, urlcolor=blue}

%% ── Colour palette (matches slides) ────────────────────────
\definecolor{iitmnavy}{RGB}{0,32,91}
\definecolor{iitmblue}{RGB}{0,90,160}
\definecolor{iitmred}{RGB}{180,20,20}
\definecolor{iitmgold}{RGB}{179,138,0}
\definecolor{notesbg}{RGB}{245,248,255}
\definecolor{tipbg}{RGB}{255,250,235}
\definecolor{warnbg}{RGB}{255,242,242}

%% ── Section style: Slide N ──────────────────────────────────
\titleformat{\section}[block]
  {\color{iitmnavy}\large\bfseries\filright}
  {Slide~\thesection.}{0.5em}{}
  [\color{iitmnavy}\titlerule]

\titlespacing{\section}{0pt}{14pt}{6pt}

%% ── Subsection: slide subtitle ──────────────────────────────
\titleformat{\subsection}[runin]
  {\color{iitmblue}\normalsize\bfseries}
  {}{0pt}{}[\ \textendash\ ]
\titlespacing{\subsection}{0pt}{4pt}{4pt}

%% ── Callout boxes ───────────────────────────────────────────
\newmdenv[
  backgroundcolor=notesbg,
  linecolor=iitmblue,
  linewidth=1.2pt,
  leftline=true, rightline=false, topline=false, bottomline=false,
  innerleftmargin=8pt, innerrightmargin=6pt,
  innertopmargin=5pt, innerbottommargin=5pt,
  skipabove=6pt, skipbelow=4pt,
]{scriptbox}

\newmdenv[
  backgroundcolor=tipbg,
  linecolor=iitmgold,
  linewidth=1.2pt,
  leftline=true, rightline=false, topline=false, bottomline=false,
  innerleftmargin=8pt, innerrightmargin=6pt,
  innertopmargin=5pt, innerbottommargin=5pt,
  skipabove=6pt, skipbelow=4pt,
]{tipbox}

\newmdenv[
  backgroundcolor=warnbg,
  linecolor=iitmred,
  linewidth=1.2pt,
  leftline=true, rightline=false, topline=false, bottomline=false,
  innerleftmargin=8pt, innerrightmargin=6pt,
  innertopmargin=5pt, innerbottommargin=5pt,
  skipabove=6pt, skipbelow=4pt,
]{warnbox}

\newcommand{\script}[1]{\begin{scriptbox}\textbf{Script:} #1\end{scriptbox}}
\newcommand{\tip}[1]{\begin{tipbox}\textbf{Tip:} #1\end{tipbox}}
\newcommand{\anticipate}[1]{\begin{warnbox}\textbf{Anticipated question:} #1\end{warnbox}}

%% ── Math macros ─────────────────────────────────────────────
\newcommand{\kibr}{k_{\mathrm{ibr}}}
\newcommand{\TMS}{\mathit{TMS}}
\newcommand{\CTI}{\mathit{CTI}}
\newcommand{\PINN}{\textsc{pinn}}
\newcommand{\WAPC}{\textsc{wapc}}

%% ── Header / footer ─────────────────────────────────────────
\pagestyle{fancy}
\fancyhf{}
\fancyhead[L]{\small\color{iitmnavy}\textbf{TSA Speaker Notes}\ ---\ Anoop V Eluvathingal}
\fancyhead[R]{\small\color{iitmgray}Self-Adaptive Model-Based Protection}
\fancyfoot[C]{\small\thepage}
\renewcommand{\headrulewidth}{0.4pt}
\renewcommand{\headrule}{\color{iitmnavy}\hrule width\headwidth height\headrulewidth}

\definecolor{iitmgray}{RGB}{90,90,90}

%% ─────────────────────────────────────────────────────────────
\begin{document}

%% ── Cover ────────────────────────────────────────────────────
\begin{center}
  {\LARGE\bfseries\color{iitmnavy}
    Thesis Submission Approval Meeting\\[6pt]
    Speaker Notes}
  \vspace{6pt}\\
  {\large Self-Adaptive Model-Based Protection for Future Power Systems}\\[4pt]
  {\normalsize Anoop V Eluvathingal\ $\cdot$\ Department of Electrical Engineering\\
    IIT Madras\ $\cdot$\ June 2026}
\end{center}

\medskip
\begin{mdframed}[backgroundcolor=notesbg, linecolor=iitmnavy, linewidth=1pt,
  innertopmargin=8pt, innerbottommargin=8pt, innerleftmargin=10pt]
\textbf{How to use these notes:}
These notes are structured as a full speaker script for each of the 22 slides.
Each section begins with the \textit{key message} — the single takeaway the
audience should retain — followed by the detailed verbal script, suggested
emphasis points, and anticipated committee questions with guidance on how to
handle them.
Total recommended talk time: \textbf{30--35 minutes} plus 10--15 minutes Q\&A.
Per-slide time targets are shown in each section header.
\end{mdframed}

\tableofcontents
\clearpage

%% ════════════════════════════════════════════════════════════
\section{Title Slide \hfill\normalfont\small($\approx$1 min)}
\label{slide:title}

\subsection*{Key message}
Introduce yourself, your thesis title, and your guide.
Establish the seriousness and scope of the problem in one sentence.

\script{%
``Good morning. I am Anoop V Eluvathingal, Roll Number EE18D203, registered for the
Doctor of Philosophy programme in the Department of Electrical Engineering at IIT Madras.
My thesis is titled \textit{`Self-Adaptive Model-Based Protection for Future Power
Systems'}, and I am working under the guidance of Dr K Shanti Swarup.

This thesis addresses a fundamental crisis that is unfolding in power grids worldwide:
the relay protection systems that have kept our networks safe for decades were designed for
synchronous generators, and they are beginning to fail as inverter-based resources take
over the grid. This work presents a complete, mathematically grounded framework to solve
that problem.''%
}

\tip{Pause briefly after ``beginning to fail'' to let the gravity of the problem register.
Make eye contact with each committee member. Keep this slide to under 60 seconds.}

%% ════════════════════════════════════════════════════════════
\section{Presentation Outline \hfill\normalfont\small($\approx$1 min)}
\label{slide:outline}

\subsection*{Key message}
The thesis is organised into four coherent parts. Set expectations clearly
so the committee can follow the arc.

\script{%
``I have organised this presentation into four parts.

Part~I covers the motivation and the SAMBP framework itself: why conventional protection
fails under high IBR penetration, what gaps currently exist in the literature, and the
architecture of the solution I am proposing.

Part~II covers five adaptive protection modules, one for each relay type studied in the
thesis: the line differential relay, generator over-current, transformer differential,
distance relay, and microgrid protection.

Part~III covers three advanced layers built on top of those modules: a real-time digital
twin using an Extended Kalman Filter, a physics-informed neural network for fault
classification, and a wide-area protection coordinator.

Part~IV concludes with the cybersecurity layer, a consolidated results comparison, and a
summary of contributions and research outcomes.

I will spend the most time on Parts~I through~III, where the novel contributions are
concentrated.''%
}

\tip{Use the left-hand column of the slide as a natural anchor for `Part I and II' and
the right-hand column for `Part III and IV'. Do not read the bullets; speak from memory.}

%% ════════════════════════════════════════════════════════════
\section{The IBR Revolution \& the Protection Crisis \hfill\normalfont\small($\approx$2.5 min)}
\label{slide:motivation}

\subsection*{Key message}
Five concrete failure mechanisms show that the IBR transition is not an incremental
change — it breaks the physical assumptions baked into every existing protection standard.

\script{%
``Let me begin with the problem. The global transition to inverter-based resources ---
solar PV, wind turbines, battery storage --- is accelerating. Under the Indian government's
2030 renewable targets, the fraction of generation from IBRs will exceed sixty percent on
many substations. I define a dimensionless parameter, $\kibr$, as the ratio of IBR
apparent admittance to the total bus admittance. A fully synchronous grid is $\kibr = 0$;
a fully IBR grid is $\kibr = 1$.

The table on the slide shows what happens to the available fault current as $\kibr$ rises.
A synchronous machine contributes six times rated current during a fault. A grid-following
inverter, by design, limits its output to 1.0 to 1.2 times rated current to protect its
semiconductor switches. At $\kibr = 0.6$, the total fault current is less than two per unit.

This creates five specific failure mechanisms for conventional protection, which I have
listed on the left. First, the differential current in an 87L relay may never exceed the
fixed slope setting when $\kibr$ is low, causing a missed trip. Second, when $\kibr$ is
high and there is CT measurement error, the relay may false-trip on load current. Third,
distance relays over-reach because IBR reactive injection shifts the apparent impedance
locus. Fourth, IBR control-mode transitions during faults produce angle excursions that
fool directional elements. Fifth, in islanded microgrids, current reverses direction
during mode transitions, making directional overcurrent protection unreliable.

The critical point is that all five failures arise from a single root cause: the relay
has a fixed setting designed for $\kibr = 0$ and no mechanism to track the true,
time-varying $\kibr$ of the network.''%
}

\anticipate{``What is the basis for the $\kibr$ definition?'' ---
$\kibr$ is the ratio of the Norton equivalent admittance of IBR sources to total bus
admittance in the positive-sequence equivalent circuit. It reduces to the
generation-mix fraction under balanced conditions and is directly measurable from PMU
phasor data.}

%% ════════════════════════════════════════════════════════════
\section{Research Gaps \& Objectives \hfill\normalfont\small($\approx$2 min)}
\label{slide:gaps}

\subsection*{Key message}
Six specific gaps in the literature, and six matching objectives that this thesis
addresses one-to-one.

\script{%
``Existing work in the literature addresses IBR protection in a piecemeal fashion:
a correction factor for differential relays here, a harmonic-blocking improvement there.
What is missing is a \textit{unified, physics-aware adaptive framework} that applies
across all relay types, updates in real time, and is provably stable.

I have identified six specific gaps. First, no relay standard or product accounts for
real-time $\kibr$ as a variable setting input. Second, the protection implications of
grid-forming IBR control modes --- as opposed to grid-following --- are largely unstudied.
Third, there is no inverse-estimation framework that is simultaneously applicable to
differential, impedance, and overcurrent relays. Fourth, digital-twin approaches proposed
in the literature lack bounded, non-blocking update laws, meaning a slow convergence can
delay the protection decision. Fifth, wide-area protection systems ignore IBR-induced
coherency breakdown, which invalidates classical swing-group assumptions. Sixth, there is
no adversarially robust ML layer integrated with the protection relay chain.

Each of these gaps maps directly to one of my six objectives, which are listed on the
right. Together, they define the scope of the SAMBP framework.''%
}

\tip{You do not need to read every objective. After listing the gaps, say ``each of these
maps directly to an objective, which I will demonstrate as we go through the chapters.''}

\anticipate{``Why is it necessary to solve all six simultaneously?'' ---
The thesis shows that solving any subset in isolation still leaves the overall protection
system vulnerable. For example, an accurate EKF oracle is useless if the ML classifier is
susceptible to adversarial attack. The strength of the SAMBP architecture is precisely its
cross-layer coherence.}

%% ════════════════════════════════════════════════════════════
\section{SAMBP Framework: Three-Tier Architecture \hfill\normalfont\small($\approx$2.5 min)}
\label{slide:architecture}

\subsection*{Key message}
SAMBP is a \emph{stack}, not a collection of independent fixes. The EKF oracle feeds
all modules; the WAPC integrates their outputs; the security layer hardens the entire chain.

\script{%
``The SAMBP framework is organised into three tiers, as shown in the diagram on this
slide.

Tier~1 is the Parameter Oracle. This is the EKF-based digital twin described in
Chapter~9. It runs continuously on PMU phasor data and maintains real-time estimates of
the three key latent parameters: $\hat{\kibr}$, the line impedance $\hat{Z}_{\rm line}$,
and the IBR admittance $\hat{Y}_{\rm ibr}$. These estimates are made available to all
modules over a low-latency local bus.

Tier~2 is the collection of adaptive protection modules. Each module --- 87L, SyncOC,
87T, distance, and microgrid --- queries the oracle at every update cycle and adjusts its
operating thresholds accordingly via what I call the bounded update law. Critically, each
module also incorporates a Stage-2 confidence gate: before issuing a trip command, the
module computes a confidence score $\gamma$ based on the uncertainty in $\hat{\kibr}$,
and only trips if $\gamma$ exceeds a configurable threshold of 0.85. This eliminates
false trips caused by transient estimation noise.

Tier~3 is the Wide-Area Protection Coordinator. It receives trip candidates from all
zone relays, maintains a real-time topology graph, and applies an IBR-aware coherency
model to determine the minimum-cost fault isolation action.

Below all three tiers sits the cybersecurity layer, which authenticates every IED message
and adversarially hardens the PINN classifier.

The key architectural principle is that the Stage-2 confidence gate decouples the speed
of the EKF oracle from the correctness of the trip decision. Even if the oracle is still
converging after a topology change, the relay will hold its last trusted estimate rather
than issuing a potentially incorrect trip.''%
}

\anticipate{``What happens if the EKF oracle fails?'' ---
Each module has a safe fallback: if the oracle has not provided an update within a
configurable timeout (default 3 cycles), the module reverts to its nominal fixed-parameter
setting and the Stage-2 gate forces the confidence score to zero, effectively disabling
adaptive operation until the oracle recovers. This is the fail-secure behaviour described
in Chapter~9.}

%% ════════════════════════════════════════════════════════════
\section{Unified Inverse-Estimation Framework \hfill\normalfont\small($\approx$2 min)}
\label{slide:math}

\subsection*{Key message}
All SAMBP modules share a single Levenberg--Marquardt inverse-estimation core.
Convergence is proved analytically; the bounded update law is the key stability ingredient.

\script{%
``Chapter~2 establishes the mathematical foundation shared by every module in the
framework.

The central problem is: given a stream of phasor measurements $z_t$ from IEDs and PMUs,
estimate the latent protection-relevant parameters $\boldsymbol{\theta} = \{\kibr,
Z_{\rm line}, Y_{\rm ibr}\}$ in real time. This is an inverse problem, and I solve it
using the Levenberg--Marquardt algorithm, which is a damped Gauss--Newton method.
The update rule is shown on the left side of the slide. The Jacobian $J$ captures how
the relay model output changes with each parameter; the damping term $\lambda I$
regularises the inverse and controls the step size.

The bounded update law is the modification I introduce. The learning rate $\alpha_t$ is
not constant; it is multiplied by a factor $(1 - \kibr^2)$ which smoothly drives the
update rate to zero as the network approaches fully IBR operation. This prevents the
estimator from taking large steps in uncertain territory and is the mechanism that enables
the Lyapunov stability proof in Theorem 2.1.

The Stage-2 confidence gate is computed as a Gaussian function of the distance between
the current estimate $\hat{\boldsymbol{\theta}}$ and the last locally-validated operating
point $\boldsymbol{\theta}^*$. The threshold $\gamma_{\rm min} = 0.85$ was chosen to give
a false-trip probability below $10^{-4}$ under CT noise levels consistent with IEC 60044.

The five methods listed in the table on the right are not independent tools. They are
layered: LM provides the parameter estimates; EKF tracks their dynamics; SPRT tests
hypotheses against those estimates; PINN provides a physics-consistent feature space;
and projected gradient enforces coordination constraints.''%
}

\anticipate{``Is LM not computationally expensive for real-time use?'' ---
The parameter vector is only 3--5 dimensional. A single LM iteration involves a 5-by-5
matrix inversion, which takes under 50 microseconds on standard IED hardware. The EKF
runs in a separate thread at 20~ms intervals; the LM correction is applied asynchronously.
Full timing budgets are detailed in Appendix~C.}

%% ════════════════════════════════════════════════════════════
\section{IBR Characterisation \& Failure Modes \hfill\normalfont\small($\approx$2 min)}
\label{slide:ibr}

\subsection*{Key message}
Chapter~3 derives the quantitative conditions for each failure mode and establishes
$\kibr$ as the single parameter that predicts them all.

\script{%
``Chapter~3 provides the analytical foundation for the failure mode analysis. I develop
a generalised IBR fault current model: the IBR contribution during a fault is $\kibr$
times the rated current, at a control angle $\phi_{\rm ctrl}$ that is set by the LVRT
strategy and is not directly observable by the relay.

The three failure modes are shown in the table. At low $\kibr$, below about 0.2, the
differential current $I_{\rm diff}$ never reaches the fixed slope boundary --- the relay
sees a fault that looks like normal load unbalance and restrains. At high $\kibr$ combined
with CT measurement error, the load current at full IBR output can exceed the slope
boundary and trigger a false trip. The third mode is impedance overreach: IBR reactive
injection adds a component to the apparent impedance that pushes the locus outside the
Zone~1 boundary even for an in-zone fault.

The insight from this analysis is that all three modes are characterised by a different
regime of $\kibr$. A relay that knows $\kibr$ can adapt its settings to remain in the
correct operating region across the entire range. This is precisely what the adaptive
slope law in Chapter~4 achieves.''%
}

\tip{Point to the table values as you describe each failure mode. Emphasise that this
is a \emph{unified} characterisation --- not three separate problems but one problem
viewed from three angles.}

%% ════════════════════════════════════════════════════════════
\section{Ch~4: SAMBP-87L --- Adaptive Line Differential (Method) \hfill\normalfont\small($\approx$2 min)}
\label{slide:87l_method}

\subsection*{Key message}
The adaptive slope law is a smooth, monotonically decreasing function of $\kibr$;
it has a closed-form Lyapunov stability proof.

\script{%
``Chapter~4 applies the SAMBP framework to the line differential relay, Type~87L.

The adaptive slope law is the core contribution. The conventional relay uses a fixed
slope $k_m = 0.30$. I replace this with a function of the estimated $\kibr$:
$k_m(\kibr) = k_{\rm base} \cdot (1-\kibr)^2 / (1+\kibr)^2 + k_{\rm floor}$.

This function has a natural interpretation. The numerator $(1-\kibr)^2$ shrinks as IBR
contribution grows, reducing the slope and keeping the relay sensitive. The denominator
$(1+\kibr)^2$ prevents the slope from ever reaching zero: as $\kibr \to 1$, $k_m$
approaches $k_{\rm floor}$, a small positive value that guards against measurement noise.

The stability proof uses the Lyapunov function $V = \|\kibr - \hat{\kibr}\|^2$.
Under bounded CT noise and the bounded update law, $\dot{V} < 0$, which guarantees that
the estimated $\kibr$ converges to a neighbourhood of the true value. The radius of that
neighbourhood is of order $\sigma_{\rm noise}$, the CT noise standard deviation.

The relay logic involves six steps. Phasors are sampled at 1~kHz; $\kibr$ is estimated
every 20~ms via the LM update; the adaptive slope is recomputed; the operate/restrain
criterion is evaluated; the Stage-2 confidence gate is checked; and the trip command is
issued only if both stages pass.''%
}

\anticipate{``How is the relay stable if $\kibr$ itself is uncertain?'' ---
The Lyapunov proof shows stability in the \emph{error} dynamics, not in the absolute
value. Even if $\hat{\kibr}$ lags the true $\kibr$ by up to one per-unit standard
deviation, the slope law still keeps the differential threshold within the safe operating
envelope. The Stage-2 gate provides an additional margin.}

%% ════════════════════════════════════════════════════════════
\section{Ch~4: SAMBP-87L --- Validated Results \hfill\normalfont\small($\approx$1.5 min)}
\label{slide:87l_results}

\subsection*{Key message}
Perfect TPR of 1.000 and zero false trips across 110 fault scenarios, including
full-IBR grid conditions where fixed relays fail in nearly one in six cases.

\script{%
``I will now present the key quantitative results for the 87L module.

The test matrix covers 110 scenarios: eleven values of $\kibr$ from 0 to 1.0, five
fault types (AG, BG, CG, BCG, ABC), two CT error levels (0\% and 5\%), and two
network loading conditions. The baseline is a fixed-slope relay with $k_m = 0.30$.

The headline result is: SAMBP-87L achieves a True Positive Rate of 1.000 and a False
Positive Rate of 0.000. The fixed relay achieves a TPR of 0.783 and a FPR of 14.8\%.
The missed trips occur entirely in the $\kibr > 0.6$ region; the false trips occur in
the $\kibr > 0.4$ plus CT error region. SAMBP eliminates both failure modes simultaneously.

The average trip time is 21 ms for SAMBP versus 22 ms for the fixed relay. The Stage-2
gate adds approximately one cycle of additional decision latency, which is acceptable
under IEC 60255-27 requirements for Class~P protection systems.

The EKF oracle achieves a $\kibr$ RMSE of 0.024~pu across all scenarios, which is
consistent with the theoretical convergence bound from Chapter~2.''%
}

\tip{The committee will focus on the TPR/FPR numbers. Be prepared to explain: (a) what
constitutes a ``scenario'' in the test matrix, (b) whether the network model is
representative of real Indian grid conditions, and (c) how the results would change under
higher CT saturation.}

\anticipate{``Is a 4-bus test network sufficient to validate the approach?'' ---
The 4-bus network was chosen to provide an analytically tractable reference for the
theoretical bounds. A 14-bus IEEE test network simulation (described briefly in
Appendix~B) shows comparable TPR gains. Hardware-in-loop validation on RTDS is planned
as future work and is listed in the thesis conclusions.}

%% ════════════════════════════════════════════════════════════
\section{Ch~5: SAMBP-SyncOC --- Generator Over-Current \hfill\normalfont\small($\approx$2 min)}
\label{slide:oc}

\subsection*{Key message}
Online adaptation of the IEC inverse-time pickup setting cuts spurious pickups by 58\%
while maintaining coordination margins across the full IBR range.

\script{%
``Chapter~5 applies the SAMBP framework to generator over-current protection, the relay
type typically designated 50/51 in IEC 60255 terminology, and which I call SyncOC
because I focus on synchronous generators operating in an IBR-rich environment.

The fundamental problem is that the classical IEC inverse-time characteristic assumes a
deterministic fault current magnitude, which is set once at the commissioning stage based
on the synchronous generator's subtransient reactance. When $\kibr$ rises, the generator's
fault contribution is depressed because the IBR sources absorb a fraction of the fault
current. The relay sees less current but its pickup setting has not changed, which can
cause a spurious pickup during a heavy load event that resembles a fault at high $\kibr$.

My solution adapts the pickup current $I_{\rm set}$ as a function of $\hat{\kibr}$, using
the same LM estimator as the 87L module. The coordination constraint --- that the primary
relay trips before the backup, with a coordination time interval $\CTI$ between them ---
is enforced at every update step via a projected gradient step that keeps the solution
within the feasible coordination set.

The results show a 58\% reduction in spurious pickups, from 38 to 16 events per 100-hour
simulation, and an $R^2$ of greater than 0.999 between the adaptive curve and the
theoretical IEC characteristic. The coordination margin improves from 0.82 to 0.94 because
the adaptive pickup tracks the actual fault current more faithfully.''%
}

\anticipate{``Does the projected gradient step always converge within one update cycle?'' ---
In all 1,200 tested operating points it converged within 3 iterations, which takes
under 0.1 ms. The theoretical convergence condition is that the gradient of the
coordination constraint is locally Lipschitz, which holds for the IEC inverse curves.}

%% ════════════════════════════════════════════════════════════
\section{Ch~6: SAMBP-87T --- Transformer Differential with SPRT \hfill\normalfont\small($\approx$1.5 min)}
\label{slide:87t}

\subsection*{Key message}
Replacing 2nd-harmonic blocking with an SPRT that has an IBR-updated likelihood model
halves the decision window and reduces IBR-induced false positives by 90\%.

\script{%
``Chapter~6 deals with transformer differential protection, relay type 87T. The
classical approach to blocking false trips caused by magnetising inrush current is to
measure the second harmonic component of the differential current. If the second harmonic
exceeds 15--20\% of the fundamental, the relay restrains. This works well for
synchronous-machine-fed transformers, but Type-4 IBR inverters produce waveforms that
may have significant second-harmonic content even during genuine internal faults, and
conversely, they can produce inrush-like fundamental distortion that lacks second harmonic
content. The 2nd-harmonic criterion therefore becomes unreliable.

I replace it with the Wald Sequential Probability Ratio Test. The SPRT accumulates a
log-likelihood ratio $\Lambda_k$ at each sample. When $\Lambda_k$ crosses an upper
threshold $A$, a trip is issued; when it crosses a lower threshold $B$, the relay
restrains. The thresholds are derived analytically from the desired Type~I and Type~II
error probabilities, both set to 1\% per IEC 60255 requirements.

The key SAMBP contribution is that the $H_1$ fault likelihood model is updated in real
time using the IBR harmonic injection spectrum from $\hat{Y}_{\rm ibr}$. This keeps the
SPRT calibrated to the actual network waveform environment.

The results: inrush rejection improves from 94.1\% to 99.6\%, internal fault TPR
improves from 96.8\% to 99.1\%, IBR-induced false positives fall from 8.7\% to 0.9\%,
and the mean decision sample count halves from 14 to 7, meaning the relay trips faster.''%
}

%% ════════════════════════════════════════════════════════════
\section{Ch~7: Adaptive Distance Protection \hfill\normalfont\small($\approx$1.5 min)}
\label{slide:dist}

\subsection*{Key message}
Zone~1 overreach drops from 34\% to 4\% of events by compensating the IBR-induced
impedance shift in real time, at less than 1 ms speed penalty.

\script{%
``Chapter~7 addresses distance protection, relay type~21. The problem here is the
impedance overreach caused by IBR reactive injection.

During a fault on a line with a high-IBR source at the remote end, the IBR injects
reactive current that adds a component $\Delta Z_{\rm ibr}$ to the apparent impedance
seen at the relay. This shifts the impedance locus radially outward on the R-X plane.
If the shift is large enough, the locus falls outside the Zone~1 boundary even though
the fault is genuinely inside Zone~1.

My solution computes the compensation term $\hat{\Delta Z}_{\rm ibr}(\hat{\kibr})$ using
the LM estimator output and adds it to the Zone~1 boundary every 10 ms. The adapted
boundary moves outward to meet the shifted locus, restoring correct Zone~1 coverage.

For severe IBR conditions where $\kibr > 0.5$, SAMBP additionally enables Zone~2
acceleration: it temporarily allows an instantaneous Zone~2 trip to prevent unacceptable
fault-clearance time.

The results: correct Zone~1 trips improve from 74\% to 96\%, overreach events fall from
17 to 2 out of 50 test faults, and underreach is eliminated entirely. The average trip
time increases by less than 1 ms due to the zone adaptation computation.''%
}

%% ════════════════════════════════════════════════════════════
\section{Ch~8: Microgrid Mode-Switching Protection \hfill\normalfont\small($\approx$1.5 min)}
\label{slide:mg}

\subsection*{Key message}
A Bayesian mode detector drives a mode-adaptive directional element, eliminating
the 22\% mis-operation rate at the grid-to-island transition boundary.

\script{%
``Chapter~8 addresses protection for microgrids that operate in both grid-connected and
islanded modes. The transition between modes is the most dangerous moment: the pre-fault
load flow used to set directional element polarity becomes invalid within one cycle of
the transition, and the fault current magnitude, direction, and impedance all change.

My approach uses a Bayesian mode detector that computes the likelihood of each mode
given the recent measurement sequence $z_t$ and the estimated $\kibr$. When the detector
switches to islanded mode with probability above 0.95, the directional element
automatically switches from the pre-fault load-flow reference to a polarity-agnostic
differential criterion with an adaptive pickup at 1.05 times rated current.

The key design choice is that the mode switch is \emph{latched}: once islanded mode is
confirmed, the relay does not revert until grid reconnection is explicitly confirmed by
a two-out-of-three GOOSE vote from neighbouring relays. This prevents chattering.

The results show a mode-detection accuracy of 98.7\%, a transition mis-operation rate
of 1.4\% versus 22\% for the fixed relay, and an islanded-fault TPR of 97.3\% versus
71\%. There are zero false island detections across 100 test scenarios.''%
}

%% ════════════════════════════════════════════════════════════
\section{Ch~9: Non-Blocking EKF Digital Twin \hfill\normalfont\small($\approx$2 min)}
\label{slide:ekf}

\subsection*{Key message}
The EKF digital twin achieves $\kibr$ RMSE of 0.024~pu with sub-2-cycle convergence
and is non-blocking: the relay never waits for the oracle.

\script{%
``Chapter~9 describes Tier~1 of the SAMBP architecture --- the EKF-based digital twin
that acts as the parameter oracle for all other modules.

The digital twin maintains a continuous state estimate $\hat{\boldsymbol{\theta}}$ using
the standard EKF prediction--correction cycle. The prediction step uses a first-order
Markov model for parameter evolution; the correction step uses the PMU voltage and current
phasors as observations. The observation function $h(\hat{\boldsymbol{\theta}})$ is the
forward model of the network derived from the nodal admittance matrix, updated with the
latest breaker status from the SCADA system.

The non-blocking property is the key architectural novelty. The EKF runs in a separate
execution thread with a dedicated real-time priority. If, after a topology change, the
EKF has not converged within 3 cycles --- 60 ms at 50 Hz --- the oracle signals a hold
condition, and each protection module continues operating with the last validated
$\hat{\boldsymbol{\theta}}$. The module does not wait. This ensures that even transient
oracle failures do not introduce protection delays.

The bounded update law prevents adversarial data injection from causing large parameter
steps. Any correction that would change $\boldsymbol{\theta}$ by more than $\delta_{\rm max}$
per cycle is clipped and flagged as a potential integrity anomaly.

The results show an RMSE of 0.024~pu for $\hat{\kibr}$, 0.018~pu for
$\hat{Z}_{\rm line}$, and 0.031~pu for $\hat{Y}_{\rm shunt}$, across all operating
conditions including 5\% PMU measurement noise. Convergence after a topology change
takes fewer than 2 cycles, and the maximum update latency is 1.8 ms.''%
}

\anticipate{``How does the EKF perform under bad data, for example a stuck PMU?'' ---
The EKF residual monitor flags measurements whose innovation exceeds three standard
deviations of the expected measurement noise. Flagged measurements are excluded from
the correction step. If all measurements from a PMU are flagged for more than 5
consecutive cycles, the oracle issues a degraded-mode signal and modules fall back to
conservative fixed settings. Full analysis is in Chapter~9, Section~4.}

%% ════════════════════════════════════════════════════════════
\section{Ch~10: Physics-Informed Fault Classifier \hfill\normalfont\small($\approx$2 min)}
\label{slide:pinn}

\subsection*{Key message}
Adding a KCL physics residual to the neural-network loss function boosts accuracy on
high-IBR faults from 64\% to 94\% --- a 46\% relative gain.

\script{%
``Chapter~10 develops the physics-informed neural network fault classifier. Pure
machine-learning approaches trained on historical fault data fail on IBR-modified
waveforms because the training data distribution is dominated by synchronous-machine
faults. Adding a physics residual to the loss function forces the network to respect
Kirchhoff's Current Law at each bus, which is the constraint that all fault types share
regardless of the source type.

The PINN loss has two terms. The first is the standard cross-entropy classification
loss on labelled fault data. The second is the KCL residual: the squared norm of
$\mathbf{K}\hat{\mathbf{V}} - \hat{\mathbf{I}}$, where $\mathbf{K}$ is the nodal
admittance matrix and $\hat{\mathbf{V}}, \hat{\mathbf{I}}$ are the voltage and current
phasors predicted by the network for the current input. The weight $\lambda_{\rm phys} = 0.3$
was chosen by cross-validation to balance the two terms.

The input feature vector includes eight elements: the three phase current magnitudes and
angles, the zero-sequence and negative-sequence components, and crucially, the oracle's
$\hat{\kibr}$ and $\phi_{\rm ctrl}$ estimates. Including the oracle outputs as features
is what allows the PINN to generalise across the full $\kibr$ range.

The results: overall accuracy rises from 81.4\% for a standard DNN to 96.1\% for the
PINN. For high-IBR conditions with $\kibr > 0.6$, the gain is from 64.2\% to 93.8\%.
Physical consistency --- meaning the predicted fault state satisfies KCL to within a
tolerance of $10^{-3}$ per unit --- rises from 51.2\% to 99.7\%.''%
}

\anticipate{``How was the PINN training dataset constructed?'' ---
The training set comprises 15,000 simulated fault scenarios generated using the
parameterised IBR fault model from Chapter~3, with $\kibr$ uniformly sampled from
$[0, 1]$, five fault types, three fault impedance levels, and three loading conditions.
Details and the dataset are available on IEEE DataPort (reference given in the thesis).}

%% ════════════════════════════════════════════════════════════
\section{Ch~11: Wide-Area Protection Coordinator \hfill\normalfont\small($\approx$2 min)}
\label{slide:wapc}

\subsection*{Key message}
Replacing electromechanical swing-group coherency with an IBR-aware coherency matrix
improves fault localisation from 71\% to 94\% and halves backup clearance time.

\script{%
``Chapter~11 describes the Wide-Area Protection Coordinator at Tier~3 of the SAMBP stack.

Classical wide-area protection systems group generators into coherent swing groups
using electromechanical swing equations. In a network with 60\% or more IBR, those
swing groups are invalid: inverters do not participate in electromechanical oscillations.
The coherency grouping therefore assigns incorrect backup priorities.

My replacement is the IBR coherency matrix $\mathbf{C}(\kibr)$, which groups buses based
on their combined IBR admittance vector similarity rather than inertia-weighted swing
angles. The matrix is computed from the oracle's $\hat{\kibr}$ values at each bus and
is updated every 100~ms, which is fast enough to track topology changes but slow enough
not to overload the communication system.

The WAPC server receives trip candidates from all zone relays over IEC 61850 GOOSE
messages. It builds a topology graph from real-time breaker status, applies the coherency
matrix to identify the fault island, and solves an Integer Linear Programme to find the
minimum-cost set of breaker operations that isolates the fault while retaining maximum
load.

The ILP has a time limit of 50~ms per decision. In all tested scenarios it solved within
35~ms. For complex topologies with more than 20 candidate breakers, a greedy warm-start
is used which finds a near-optimal solution within 10~ms.

The results: fault localisation accuracy rises from 71.2\% to 94.5\%, unnecessary trips
fall from 18.4\% to 4.1\%, load preserved at the end of the isolation sequence rises
from 63\% to 87.6\%, and backup fault clearance time is halved from 420 ms to 210 ms.''%
}

\anticipate{``What is the communication latency assumed?'' ---
The model assumes a worst-case one-way GOOSE latency of 4~ms per hop, consistent with
IEC 61850 Performance Class P2. Sensitivity analysis in Chapter~11, Section~5, shows
that the fault-localisation accuracy remains above 90\% up to a latency of 12~ms, beyond
which the ILP topology snapshot becomes stale.}

%% ════════════════════════════════════════════════════════════
\section{Ch~12: Cybersecurity Architecture \hfill\normalfont\small($\approx$1.5 min)}
\label{slide:sec}

\subsection*{Key message}
HMAC authentication and adversarial PINN training reduce the composite risk index by
54.9\% with less than 0.3 ms overhead — within IEC 62351 latency budgets.

\script{%
``Chapter~12 addresses the cybersecurity of the SAMBP architecture. As the system
relies on real-time PMU data and IED messages, it presents two attack surfaces that
I identify and defend against.

The first threat is data integrity attack. An adversary who can modify PMU measurements
in transit can inject false $\hat{\kibr}$ values and induce the protection system to
mal-operate. The defence is HMAC-SHA256 authentication on every IED message, using keys
managed by the IEC 62351-8 key management infrastructure. A relay that receives an
unauthenticated or tampered message discards it and holds the last trusted parameter
estimate. The authentication overhead is 0.3 ms per message, which is within the
IEC 62351 latency budget for Class~P2 systems.

The second threat is adversarial machine learning attack. An adversary who knows the
PINN architecture can craft a perturbation $\delta\mathbf{x}$ with magnitude less than
5\% of the input range that causes the classifier to mis-identify a fault. Without
defence, this drops the PINN accuracy from 96\% to under 52\%.

The defence is adversarial training using Fast Gradient Sign Method and Projected
Gradient Descent examples. The retrained network achieves a certified robustness radius
of 0.04~pu, meaning no perturbation within that radius can change the classification.
Accuracy under attack recovers to 91.4\%.

The composite risk index --- a weighted sum of attack success rate, impact severity, and
relay latency --- falls from 1.000 to 0.451, a 54.9\% reduction.''%
}

%% ════════════════════════════════════════════════════════════
\section{Consolidated Performance Across All Modules \hfill\normalfont\small($\approx$1.5 min)}
\label{slide:consolidated}

\subsection*{Key message}
Every SAMBP module outperforms its fixed-parameter baseline, and the gains are
consistent across the full IBR penetration range $\kibr \in [0, 1]$.

\script{%
``Let me now bring together the results from all chapters into a single comparison table.

Looking across the rows, every SAMBP module shows a clear improvement over its
fixed-parameter baseline. The 87L relay goes from 78\% TPR to 100\% and eliminates
all false trips. The SyncOC relay cuts spurious pickups by 58\%. The 87T relay reduces
IBR-induced false positives by 90\%. The distance relay improves correct Zone~1
discrimination from 74\% to 96\%. The microgrid relay cuts transition mis-operations
from 22\% to 1.4\%. The PINN classifier improves from 81\% to 96\% overall accuracy.
The WAPC improves fault localisation from 71\% to 94\%. And the cybersecurity layer
cuts the composite risk index nearly in half.

A key aspect of these results is that they are measured under the \emph{same} test
conditions as the baselines, with no tuning advantage given to SAMBP. The only
difference is the adaptive parameter update. This validates the core claim of the thesis:
that real-time knowledge of $\kibr$ is sufficient to repair the protection failures caused
by the IBR transition, without requiring changes to relay hardware or communication
infrastructure.''%
}

\tip{Spend about 30 seconds on this slide — the committee will read the table themselves.
Focus on the narrative: one cause (unknown $\kibr$), one solution (adaptive estimation),
consistent gains across all relay types.}

%% ════════════════════════════════════════════════════════════
\section{Conclusions \& Contributions \hfill\normalfont\small($\approx$2 min)}
\label{slide:conclusions}

\subsection*{Key message}
Six principal contributions, each novel by a specific technical criterion; four directions
for future work that connect directly to real-world deployment.

\script{%
``Let me summarise the principal contributions of this thesis.

First, the SAMBP framework itself. To my knowledge, this is the first unified
inverse-estimation architecture that applies a single parameter oracle to multiple relay
types simultaneously. Prior work has proposed adaptive corrections for individual relay
types in isolation; this thesis provides the theoretical and implementation framework that
unifies them.

Second, the adaptive 87L result. A TPR of 1.000 with zero false trips across the full
IBR range is a result that no prior adaptive differential relay paper has reported. The
key enabling technology is the Stage-2 confidence gate, which I believe is also a novel
contribution to IED trip-chain design.

Third, the non-blocking EKF digital twin. The non-blocking property --- where the relay
never waits for the oracle --- is a specific design decision that, again to my knowledge,
has not been formalised in prior digital-twin protection literature.

Fourth, the PINN fault classifier. The use of a KCL physics residual as a loss-function
term is established in the scientific ML literature, but this is its first application to
power-system protection classification under IBR conditions.

Fifth, the IBR-aware WAPC. The substitution of the IBR coherency matrix for
electromechanical swing groups is a conceptually clean contribution that directly addresses
why wide-area protection systems are degrading in practice.

Sixth, the integrated cybersecurity layer. Combining HMAC authentication with adversarial
PINN training as a unified security architecture for a protection stack is a novel design.

For future work, the three most important directions are: RTDS hardware-in-loop
validation, field deployment on a representative portion of the TANGEDCO network in Tamil
Nadu, and extension of the IBR model to grid-forming inverters with virtual inertia, which
introduce dynamics not captured by the current $\kibr$ formulation.''%
}

%% ════════════════════════════════════════════════════════════
\section{Research Outcomes \hfill\normalfont\small($\approx$1 min)}
\label{slide:outcomes}

\subsection*{Key message}
Four journal papers, two conference papers, one patent, one codebase, two datasets —
all derived directly from the thesis chapters.

\script{%
``The research in this thesis has produced the following outcomes.

Four journal papers are under review. The 87L work is submitted to IEEE Transactions on
Power Delivery. The SyncOC work is submitted to IET Generation, Transmission and
Distribution. The PINN classifier paper is under review at IEEE Transactions on Neural
Networks and Learning Systems. The WAPC paper is being prepared for Electric Power Systems
Research.

Two conference papers have been published: the EKF digital-twin work was presented at
IEEE PES ISGT-Europe 2025, and the adversarial robustness paper was presented at
IEEE TENCON 2025.

A patent application has been filed through IIT Madras on the adaptive relay system with
real-time IBR parameter estimation.

The SAMBP simulation framework is approximately 4,200 lines of Python and MATLAB code
and is available on request. Two fault-signature datasets for IBR-grid conditions have
been uploaded to IEEE DataPort and are publicly accessible.''%
}

\tip{Keep this brief. If a committee member asks for a specific paper, say the preprint
is available and offer to share it after the meeting.}

%% ════════════════════════════════════════════════════════════
\section{Thesis Content \& Chapter--Publication Mapping \hfill\normalfont\small($\approx$1 min)}
\label{slide:content}

\subsection*{Key message}
The thesis is 13 chapters plus four appendices, approximately 204 pages, with clear
traceability from chapters to publications.

\script{%
``This slide shows the complete structure of the thesis. All thirteen chapters and four
appendices are listed with estimated page counts and their associated publication outputs.

The thesis begins with an introduction and mathematical foundations in Chapters~1 and~2.
Chapters~3 through~12 are the research chapters, each corresponding to one module or
cross-cutting concern of the SAMBP framework. Chapter~13 contains the conclusions and
future-work directions.

The four appendices contain the proof details from Chapter~2, the full system data for
all test networks, the pseudocode listings for all algorithms, and additional data tables
for the SyncOC coordination study.

The total estimated length is approximately 204 pages, which is consistent with IIT Madras
PhD dissertation norms for Electrical Engineering.''%
}

%% ════════════════════════════════════════════════════════════
\section{Thank You / Questions \hfill\normalfont\small(Q\&A)}
\label{slide:thanks}

\subsection*{Key message}
Leave the committee with the three headline numbers: TPR=1.000, 96.1\% accuracy,
$-$54.9\% risk reduction.

\script{%
``Thank you for your attention. To summarise in three numbers: the adaptive line
differential relay achieves a true positive rate of 1.000 --- no missed trips, no false
trips. The physics-informed fault classifier achieves 96.1\% overall accuracy with 93.8\%
on high-IBR scenarios where conventional deep networks fall to 64\%. And the cybersecurity
architecture reduces the composite risk index by 54.9\% with less than 0.3 milliseconds
of authentication overhead.

I am happy to take your questions.''%
}

\tip{Stand near the screen and maintain eye contact with the committee. Do not start
listing additional results unless a committee member specifically asks. If you do not
know an answer, say so directly: ``That is a good question; I have not analysed that
case and it would be an interesting direction.'' Confidence under uncertainty is more
reassuring to a doctoral committee than a speculative answer.}

\bigskip
\begin{mdframed}[backgroundcolor=notesbg, linecolor=iitmnavy, linewidth=1pt,
  innertopmargin=8pt, innerbottommargin=8pt, innerleftmargin=10pt]
\textbf{Frequently anticipated Q\&A topics (prepare these in advance):}
\begin{enumerate}[leftmargin=1.5em, itemsep=4pt]
  \item How does the method scale to a 400-bus transmission network? \\
        \textit{The LM estimator is per-feeder; the EKF operates on the local
        equivalent circuit. The WAPC ILP complexity scales with the number of
        candidate isolation breakers, not total bus count. Full scalability analysis
        is in Chapter~11, Section~6.}
  \item What is the minimum PMU density required? \\
        \textit{One PMU per protected feeder terminal is sufficient for the EKF.
        The WAPC requires at least one PMU per coherency group. Sensitivity to
        PMU dropout is analysed in Chapter~9.}
  \item How does SAMBP interact with existing distance relay firmware? \\
        \textit{The adaptive zone boundaries are output as IEC 61850 XSWI setting
        messages to the existing IED. The relay firmware is unchanged. This is a
        key practical design constraint of the architecture.}
  \item Has the method been tested on real utility data? \\
        \textit{Not yet. All results are based on high-fidelity simulation using
        validated IBR models. RTDS validation and TANGEDCO field trials are
        listed as future work.}
  \item What is the cost of deploying the SAMBP oracle on existing RTU hardware? \\
        \textit{The EKF oracle requires approximately 50 MIPS of compute, consistent
        with mid-range IED processors (e.g., Siemens SIPROTEC~5 class). The codebase
        is written in C-compatible Python and has been profiled on ARM Cortex-A53.}
\end{enumerate}
\end{mdframed}

\end{document}
